\documentclass[11pt,a4paper]{article}

\usepackage[utf8]{inputenc}
\usepackage[T1]{fontenc}
\usepackage{lmodern}
\usepackage[a4paper,margin=2.5cm]{geometry}
\usepackage{hyperref}
\usepackage{booktabs}
\usepackage{longtable}
\usepackage{listings}
\usepackage{xcolor}

\hypersetup{colorlinks=true,linkcolor=blue,urlcolor=blue}
\definecolor{codegray}{gray}{0.95}
\lstset{basicstyle=\ttfamily\small,backgroundcolor=\color{codegray},frame=single,breaklines=true,columns=fullflexible}

\title{HeatingSystemController\\Felhasználói útmutató}
\author{AlfaOmega projekt}
\date{\today}

\begin{document}
\maketitle
\tableofcontents
\newpage

\section{Rövid áttekintés}

A HeatingSystemController egy WiFi-n keresztül elérhető fűtésvezérlő rendszer. Kezeli a kazánt, a puffert, a radiátoros/fűtési köröket, a DS18B20 hőmérséklet-szenzorokat és a PWM/ON/OFF kimeneteket.

A rendszer webes felületről, Postmanből vagy más HTTP kliensből vezérelhető.

A firmware a soros kimenet mellett szöveges alkalmazáslogot is ment. A logfájlok összesített mérete legfeljebb 10 KB lehet; amikor az új bejegyzés már nem férne el, a legrégebbi logfájlok automatikusan törlődnek.

A fájlírás nem minden logüzenetnél történik. A rendszer először memóriapufferben gyűjti a bejegyzéseket, majd a puffer megtelésekor vagy néhány másodpercenként egyben írja őket a fájlba. Emiatt egy friss logbejegyzés rövid ideig csak a soros kimeneten látható, amíg a puffer ki nem ürül.

\section{Elérés}

A vezérlő IP-címét a hálózati beállítások vagy a soros bootnapló alapján lehet megtudni. Az API alapértelmezett portja 80.

Példa:

\begin{lstlisting}
http://192.168.0.107
\end{lstlisting}

Ha a vezérlő hostname-je elérhető a helyi hálózaton, az IP-cím helyett használható a hostname is.

\section{Aktuális állapot lekérése}

Postman beállítása:

\begin{itemize}
  \item Method: \texttt{GET}
  \item URL: \texttt{http://ESZKOZ-IP/state}
  \item Body: nincs
  \item Header: \texttt{Accept: application/json}
\end{itemize}

A válaszban megtalálhatók a kimenetek, szenzorértékek és modulállapotok.

\begin{lstlisting}[language=json]
{
  "1": 0,
  "3": 50,
  "28000B96A00000000029": 55.4,
  "Puffer950l": {
    "state": "CHARGE",
    "HT": "12.3"
  }
}
\end{lstlisting}

\section{Teszt hőmérsékletek beállítása}

Ez a végpont a tesztkörnyezetben használható szenzorértékek beállítására.

\begin{lstlisting}
POST http://ESZKOZ-IP/setTemp
Content-Type: application/json
\end{lstlisting}

Példa Body:

\begin{lstlisting}[language=json]
{
  "28000B96A00000000029": 55.0,
  "2814576B00000000059": 48.0,
  "2881506B00000000052": 60.0
}
\end{lstlisting}

Sikeres válasz:

\begin{lstlisting}[language=json]
{"status":"success"}
\end{lstlisting}

\section{Fűtési igény beállítása}

\begin{lstlisting}
GET http://ESZKOZ-IP/heatingRequest?name=PadloKozos&setTemp=true
\end{lstlisting}

A fűtési igény kikapcsolása:

\begin{lstlisting}
GET http://ESZKOZ-IP/heatingRequest?name=PadloKozos&setTemp=false
\end{lstlisting}

A \texttt{name} értékének szerepelnie kell a konfigurált fűtési elemek között.

\section{Kézi pumpavezérlés}

\begin{lstlisting}
POST http://ESZKOZ-IP/manualControlPumps
Content-Type: application/json
\end{lstlisting}

Példa: a 3-as kimenet 50\%-ra kézi módba állítása, a 4-es kimenet visszaadása automatikus módba:

\begin{lstlisting}[language=json]
{
  "3": {"manual": true, "value": 50},
  "4": {"manual": false}
}
\end{lstlisting}

A PWM érték tartománya 0 és 100. Kézi módban az automatikus vezérlés nem módosítja az adott kimenetet.

\section{Állapotnapló használata}

A vezérlő napi bináris fájlokat tárol a belső SPIFFS fájlrendszerben:

\begin{lstlisting}
/state_YYYY-MM-DD.bin
\end{lstlisting}

A fájlokat nem kell közvetlenül megnyitni vagy dekódolni. Az API JSON formátumban visszaadja a tartalmat.

\subsection{Elérhető naplófájlok}

\begin{lstlisting}
GET http://ESZKOZ-IP/state-history/files
\end{lstlisting}

Példa válasz:

\begin{lstlisting}[language=json]
{
  "currentDate": "2026-11-06",
  "files": [
    {"name":"state_2026-11-05.bin","size":6213}
  ]
}
\end{lstlisting}

A \texttt{currentDate} az eszköz aktuális, NTP által beállított dátuma. Ha a frontend nem kapott külön dátumválasztást, ezt a dátumot használja alapértelmezettként.

\subsection{Egy nap adatainak lekérése}

\begin{lstlisting}
GET http://ESZKOZ-IP/state-history?date=2026-11-05
\end{lstlisting}

A válasz csak a változásokat tartalmazza:

\begin{lstlisting}[language=json]
[
  {
    "timestamp": 1793913984,
    "changes": {
      "io/1": 0,
      "sensor/2814576B00000000059": 48.5,
      "state/Puffer950l/state": "CHARGE"
    }
  }
]
\end{lstlisting}

Fontos, hogy a későbbi rekordok nem tartalmazzák a teljes állapotot. A frontendnek az első rekordból kell létrehoznia az állapotot, majd minden további rekord \texttt{changes} mezőjét rá kell másolnia az előző állapotra.

\subsection{Naplófájl törlése}

\begin{lstlisting}
DELETE http://ESZKOZ-IP/state-history?date=2026-11-05
\end{lstlisting}

Sikeres válasz:

\begin{lstlisting}[language=json]
{"status":"deleted"}
\end{lstlisting}

A törlés előtt mindig ellenőrizd a dátumot, mert a művelet nem visszavonható.

\section{Eszközinformáció és szabad tárhely}

\begin{lstlisting}
GET http://ESZKOZ-IP/deviceInfo
\end{lstlisting}

A válasz \texttt{spiffs} része megmutatja a teljes, használt és szabad fájlrendszer-kapacitást:

\begin{lstlisting}[language=json]
"spiffs": {
  "total": 1048576,
  "used": 245760,
  "free": 802816
}
\end{lstlisting}

\section{Konfiguráció lekérése}

\begin{lstlisting}
GET http://ESZKOZ-IP/configuration
\end{lstlisting}

Ez visszaadja a jelenlegi \texttt{config.json} tartalmát. A konfigurációban lehet megadni a WiFi-t, a kazánt, a puffert, a radiátorköröket, a szenzorokat, a pumpákat és a PID paramétereket.

\section{Konfiguráció feltöltése}

A konfiguráció feltöltése a következő végponton történik:

\begin{lstlisting}
POST http://ESZKOZ-IP/configuration
Content-Type: multipart/form-data
\end{lstlisting}

Postmanben a Body alatt a \texttt{form-data} típust válaszd, majd a fájl mezőben töltsd fel a valid konfigurációt \texttt{config.json} néven.

Sikeres feltöltés után a vezérlőt újra kell indítani. A rendszer az új konfigurációt validálja, és hibás konfiguráció esetén a korábbi konfigurációt megtartja.

\section{Biztonsági megjegyzések}

A konfigurációs fájl WiFi-jelszót is tartalmazhat, ezért valódi jelszavakat ne tölts fel forráskódba és ne ossz meg dokumentációban.

A jelenlegi újraindítási végpont fejlesztői hitelesítést használ:

\begin{lstlisting}
GET /reboot?user=admin&password=password
\end{lstlisting}

Éles hálózatban ezt a hozzáférést korlátozni vagy módosítani kell.

\section{Gyakori hibák}

\begin{longtable}{p{0.33\textwidth}p{0.55\textwidth}}
\toprule
Hiba & Teendő \\
\midrule
\texttt{state log not found} & Ellenőrizd a dátumot a \texttt{/state-history/files} végponttal. \\
Üres \texttt{[]} válasz & Lehet, hogy az adott naphoz még nincs minta, vagy a firmware nem futott elég ideig a logger első mentéséig. \\
HTTP 429 & Az API limiter aktív; várj a \texttt{Retry-After} fejlécben jelzett időt. \\
HTTP 400 konfiguráció feltöltéskor & A JSON szerkezete vagy egy szenzor/PID hivatkozása hibás. \\
Feltöltési kapcsolat megszakad & Ellenőrizd az USB-kábelt, a soros portot, a tápellátást és az ESP8266 FLASH módot. \\
\bottomrule
\end{longtable}

\end{document}
